\documentclass[12pt]{article}
\usepackage[T1]{fontenc}
\usepackage[utf8]{inputenc}
\usepackage{lmodern}
\usepackage{textcomp}
\usepackage{enumitem}
\usepackage{geometry}
\geometry{margin=1in}
\begin{document}
\begin{center}
Answer Key - Biology - Graduate Level Assessment 7 \
\small (Answers are ordered exactly as in the question paper.)
\end{center}

\section*{Multiple Choice}
\begin{enumerate}[label=\arabic*.]
\item \textbf{Answer:} A
\\ \textit{Options:} A) Temperature, humidity, and air currents; B) Plant species, root depth, and wind speed; C) Soil fertility, irrigation methods, and cloud cover; D) Water salinity, leaf color, and time of day; E) Leaf surface area, stem length, and moon phase
\\ \textit{Note:} Temperature, humidity, and air currents influence the rate of transpiration by affecting the evaporation of water from plant surfaces, with higher temperatures increasing evaporation rates, lower humidity reducing the moisture gradient between the leaf and the atmosphere, and air currents enhancing the removal of water vapor from around the leaves.
\item \textbf{Answer:} B
\\ \textit{Options:} A) Frame shift mutations and nucleotide repeats; B) Base substitutions and small insertions and deletions; C) Base additions and large insertions; D) Silent mutations and missense mutations; E) Duplication mutations and triplet repeat expansions
\\ \textit{Note:} The correct answer identifies the two main categories of point mutations, which are base substitutions that alter a single nucleotide and small insertions or deletions that change the length of the DNA sequence, both of which can significantly impact gene function and protein synthesis.
\item \textbf{Answer:} C
\\ \textit{Options:} A) It is created by the direct action of natural selection.; B) It tends to be reduced by the processes involved when diploid organisms produce gametes.; C) It must be present in a population before natural selection can act upon the population.; D) It is always increased by the processes involved when diploid organisms produce gametes.; E) It is not affected by the environmental changes.
\\ \textit{Note:} Genetic variation is essential for natural selection because it provides the raw material for evolutionary change, allowing certain traits to be favored over others in response to environmental pressures.
\item \textbf{Answer:} C
\\ \textit{Options:} A) Aristotle, Darwin, and Lamarck; B) Aristotle, Lyell, and Darwin; C) Aristotle, Linnaeus, and Cuvier; D) Aristotle, Cuvier, and Lamarck; E) Linnaeus, Darwin, and Lamarck
\\ \textit{Note:} Aristotle, Linnaeus, and Cuvier all upheld the belief in the fixity of species, asserting that species were created as distinct entities and did not undergo significant evolutionary changes over time.
\item \textbf{Answer:} A
\\ \textit{Options:} A) Earthworms use gills, frogs use both gills and lungs, fish use lungs, and humans use skin.; B) Earthworms and humans use lungs, frogs and fish use gills.; C) All use gills for obtaining oxygen.; D) Earthworms and humans use gills, frogs use skin, and fish use lungs.; E) Earthworms use lungs, frogs use gills, fish use skin, and humans use gills.
\\ \textit{Note:} The correct answer highlights the diverse respiratory adaptations among these organisms, where earthworms rely on cutaneous respiration through their skin, frogs utilize both cutaneous respiration and lungs during different life stages, fish predominantly extract oxygen using gills, and humans exclusively utilize lungs for breathing, showcasing the evolutionary variations in oxygen acquisition strategies across different species.
\end{enumerate}

\section*{True / False}
\begin{enumerate}[label=\arabic*.]
\setcounter{enumi}{5}
\item \textbf{Answer:} True
\\ \textit{Options:} A) True; B) False
\\ \textit{Note:} This Google-based search was useful to identify an appropriate diagnosis in complex immunological and allergic cases. Computing skills may help to get better results.
\item \textbf{Answer:} False
\\ \textit{Options:} A) True; B) False
\\ \textit{Note:} Numerous studies have found that taxation is one of the most effective policy instruments for tobacco control. However, these findings come from countries that have market economies where market forces determine prices and influence how cigarette taxes are passed to the consumers in retail prices. China's tobacco industry is not a market economy; therefore, non-market forces and the current Chinese tobacco monopoly system determine cigarette prices. The result is that tax increases do not necessarily get passed on to the retail price.
\item \textbf{Answer:} True
\\ \textit{Options:} A) True; B) False
\\ \textit{Note:} Head and neck-specific QOL measures are necessary and should include domains that reflect ES, SC, and AP.
\item \textbf{Answer:} True
\\ \textit{Options:} A) True; B) False
\\ \textit{Note:} Findings suggest that traffic law reforms in order to have an effect on both traffic fatality and injury rates reduction require changes in police enforcement practices. Last, this case also illustrates how the diffusion of successful road safety practices globally promoted by WHO and World Bank can be an important influence for enhancing national road safety practices.
\item \textbf{Answer:} True
\\ \textit{Options:} A) True; B) False
\\ \textit{Note:} In the present study, FT and androstenedione were statistically significantly correlated with sexual desire in the total cohort of women. ADT-G did not correlate more strongly than circulating androgens with sexual desire and is therefore not superior to measuring circulating androgens by mass spectrometry.
\end{enumerate}

\section*{Long Form Response}
\begin{enumerate}[label=\arabic*.]
\setcounter{enumi}{10}
\item \textbf{Answer:} Radioactive isotopes play a crucial role in biological research, particularly in tracing metabolic pathways within living cells. By incorporating these isotopes into specific biomolecules, researchers can track their movement and transformation, providing insights into metabolic processes. For example, when isotopes are used, they can change the color of cells, thus enhancing visibility and facilitating detailed observation under a microscope. This ability to visualize metabolic activity allows for more precise experimental observations and a deeper understanding of cellular functions. Overall, the use of radioactive isotopes not only aids in tracing metabolic pathways but also improves the clarity and accuracy of biological research outcomes.
\\ \textit{Note:} The answer correctly emphasizes that radioactive isotopes enable researchers to trace and visualize metabolic pathways in living cells by incorporating them into biomolecules, thereby enhancing experimental observations and understanding of metabolic processes.
\item \textbf{Answer:} Ethylene plays a crucial role in the ripening of fruits by acting as a plant hormone that regulates various physiological processes, including the transition from immature to mature fruit. In the early 1900 s, there was a prevalent misconception that the heat generated from kerosene stoves was solely responsible for this ripening process. However, it was actually the incomplete combustion products of kerosene, particularly ethylene, that facilitated the ripening of fruits. This understanding highlights the importance of ethylene as a natural signaling molecule in fruit development, challenging the earlier views that prioritized thermal effects. By recognizing ethylene's role, we can better appreciate the biochemical mechanisms underlying fruit ripening and the historical context of agricultural practices.
\\ \textit{Note:} correct because it accurately identifies ethylene as the key hormone that triggers fruit ripening and clarifies that the misconception about kerosene stoves pertains to the role of ethylene produced from incomplete combustion rather than heat alone.
\end{enumerate}

\vfill
\noindent\textit{End of Answer Key}
\end{document}